\section{方法}

本節依最終系統說明：基線模型、長文本 encoder、對比式 supervised fine-tuning、動態負樣本選取、法律領域繼續預訓練、scope filtering、learning-to-rank，以及 post-processing。下文沿用 DAPT、SFT 與 LTR 三個縮寫。

圖~\ref{fig:task1-system-pipeline} 先總覽本文最終採用的 Task~1 系統流程，區分訓練/驗證前處理、模型訓練，以及測試階段的預測與輸出。

\begin{figure*}[t]
  \centering
  \includegraphics[width=\textwidth]{../shared/figures/task1_system_pipeline_zh.pdf}
  \caption{本文最終使用的 Task~1 系統流程總覽。圖中整理了訓練/驗證前處理資產、包含動態負樣本選取的訓練迴圈，以及測試階段結合 lexical 打分、dense 打分、LightGBM 重排與最後 post-processing 的完整路徑。}
  \Description{Task 1 系統流程圖，分成訓練驗證前處理、模型訓練與測試輸出三個區塊。流程包含原始訓練案例清理與切分、BM25 與年份範圍建立、對比式訓練資料與動態負樣本驅動的對比式 supervised fine-tuning 和 learning-to-rank 訓練，以及測試資料經 lexical 打分、dense 打分、learning-to-rank 重排與最後後處理後輸出答案。}
  \label{fig:task1-system-pipeline}
\end{figure*}

\subsection{BM25 與 SAILER 基線}

我們保留 BM25 與 SAILER 作為主要基線。BM25 代表 lexical retrieval 的強基線；SAILER 則對應 THUIR@COLIEE 2023 提出的 legal encoder \cite{li2023thuir}。在本文中，SAILER 不再重新 fine-tune，而是直接使用既有模型進行嵌入與相似度計算\footnote{\url{https://huggingface.co/CSHaitao/SAILER_en_finetune}}。這個設定其實對 SAILER 偏有利，因為其公開 fine-tuning 模型很可能已經看過與本文 valid split 重疊、或至少高度相近的資料分布。

本文的 lexical features 主要來自 BM25 與 query likelihood model（QLD）。對 query $q$ 與文件 $d$，BM25 分數為
\begin{equation}
\mathrm{BM25}(q,d)=\sum_{t \in q}\mathrm{IDF}(t)\cdot
\frac{f(t,d)(k_1+1)}{f(t,d)+k_1\left(1-b+b\frac{|d|}{\mathrm{avgdl}}\right)},
\end{equation}
其中 $f(t,d)$ 為詞項 $t$ 在文件中的出現次數，$|d|$ 為文件長度，$\mathrm{avgdl}$ 為語料平均文件長度。QLD 採用 Dirichlet smoothing：
\begin{equation}
\mathrm{QLD}(q,d)=\sum_{t \in q} c(t,q)\log \frac{c(t,d)+\mu p(t \mid C)}{|d|+\mu},
\end{equation}
其中 $c(t,q)$ 與 $c(t,d)$ 為 query 與文件中的詞項次數，$p(t \mid C)$ 為集合語言模型。

\subsection{ModernBERT 長文本編碼}

為檢驗長文本截斷是否為主要瓶頸，我們採用 Answer.AI 釋出的 ModernBERT-base \cite{warner2024modernbert}\footnote{\url{https://huggingface.co/answerdotai/ModernBERT-base}} 作為基礎 encoder。ModernBERT 原生支援 8192 tokens，理論上比 SAILER 更適合處理長篇判決書；但若在對比式訓練時同時輸入 1 篇 query、1 篇 positive 與 15 篇 negatives，等同於單一步驟要編碼 17 篇長文，在單張 RTX 4080 SUPER 16GB 上會直接發生記憶體不足。因此，本文正式實驗採用 4096 tokens。

對於每篇判決書，我們使用 encoder 開頭的 `[CLS]` 向量表示作為 dense retrieval 嵌入，並在實作中再經 projector 與 L2 normalization；本文所有 dense 排名皆以 cosine similarity 計算。

\subsection{對比式 SFT 與靜態 Hard Negatives}

我們將法律判決檢索建模為 query--positive--negative 的對比式學習問題。每個對比式樣本包含 1 篇 query case、1 篇 noticed case 與 15 篇負樣本。設第 $i$ 個樣本的 query 為 $q_i$、正樣本為 $p_i$、負樣本集合為 $\mathcal{N}_i$、相似度函數為 $s(\cdot,\cdot)$、溫度為 $\tau$，則其對比式 SFT 損失為
\begin{equation}
\mathcal{L}_i=
-\log
\frac{\exp\left(s(q_i,p_i)/\tau\right)}
{\exp\left(s(q_i,p_i)/\tau\right)+\sum_{n \in \mathcal{N}_i}\exp\left(s(q_i,n)/\tau\right)}.
\end{equation}
雖然每個相似度分數是逐對計算，但這個損失會把 1 個 positive 與 15 個 negatives 一起放進同一個 softmax 中共同正規化，因此本文的 dual encoder 對比式訓練在性質上更接近 listwise，而不是傳統只比較單一正負對的 pairwise 排序。

早期實驗依照類似 THUIR 的做法，先用 BM25 top-100 建立 hard-negative pool，再於每個訓練步驟抽取 15 篇負樣本。這個設計實作簡單，但負樣本難度固定，且完全由 lexical ranking 定義。

\subsection{動態負樣本選取}

本文最重要的方法改動，是將靜態 hard negatives 改為動態負樣本選取。在每個 epoch 開始時，先用當前模型計算 query 與 candidate 的相似度，再依據分數決定負樣本被抽中的機率；分數愈高，代表模型愈容易與其混淆，該文件就愈可能被視為 hard negative。若第 $e$ 個 epoch 中 query $q_i$ 的候選負樣本池為 $\mathcal{C}_i$，則抽樣機率為
\begin{equation}
\pi_e(n \mid q_i)=
\frac{\exp\left(\alpha s_e(q_i,n)\right)}
{\sum_{m \in \mathcal{C}_i}\exp\left(\alpha s_e(q_i,m)\right)},
\quad n \in \mathcal{C}_i,
\end{equation}
其中 $s_e$ 為當前 epoch 的相似度分數，$\alpha$ 控制分佈尖銳度。本文實作中 \texttt{sampling\_temperature=1.0}，因此 $\alpha=1$。這讓負樣本難度會隨模型狀態同步更新，而不是永遠停留在 BM25 所定義的 lexical 相似層次。

\subsection{法律領域繼續預訓練}

除了 task-specific SFT 之外，我們也先在美國判決書語料上對 ModernBERT 進行 DAPT。這個階段使用 Common Pile CaseLaw Access Project subset，讓 encoder 先適應法律文本中的長句、引文格式與程序語言。受限於目前 GPU 條件，這個 DAPT 階段約訓練 60 小時，仍只覆蓋 0.384 epoch，因此本文將它視為「仍未訓練飽和」但有幫助的法律領域適應步驟。

\subsection{年份範圍約束與 scope filtering}

本文在資料建構與最終推論兩個階段都使用 year-based scope filtering。法律判決在引用既有案例時，只會引用過去已存在的判例，而不會引用未來判決。因此，對每個 query，我們先抽取其判決年份，再只讓法律上合理的舊案進入負樣本候選池與最終推論候選集合。需要強調的是，目前實作不會因 scope 規則而刪除 gold positives；若某個 noticed case 因規則式年份抽取而落在 scope 外，它仍會保留為正樣本。

在訓練時，scope filtering 能避免把未來判決混入負樣本；在推論時，我們先離線建立 query 專屬的候選範圍索引，使每個 query 只需對自身 scope 中的候選計分，而不必與全庫逐一比較，因此同時改善訓練訊號品質與計算效率。本文共用的 scope index 採用「candidate year 不晚於 query year + 1」的規則，以吸收年份抽取噪音。

\subsection{Learning-to-Rank (LTR)}

延續 THUIR 在 WSDM Cup 2023、網頁搜尋排序與 NTCIR-16 WWW-4 任務上的融合思路 \cite{chen2023wsdmcup,li2023websearch,yang2023ntcir}，本文將 dense 與 lexical 訊號整合為最終排序分數。對每個 query--candidate pair，我們建立 28 個數值特徵，可概分為七類：lexical scores、dense scores、rank positions、長度統計、placeholder 統計、年份差，以及 1--3 個 chunks 聚合 embeddings 的相似度；\texttt{query\_id} 與 \texttt{candidate\_id} 僅作為識別欄，不會進入 ranker 訓練。

排序模型採用 LightGBM 的 \texttt{LGBMRanker}，其訓練目標為 \texttt{lambdarank}，評估指標為 \texttt{ndcg}。模型會在 validation 上做 early stopping，並以 validation 表現最佳的 ranker 對 test set 推論。

\subsection{Post-processing}

在 LTR 輸出 query-wise reranked candidates 後，本文再加入兩類 post-processing。第一類是 common legal filters，包括 query-specific 年份範圍過濾與自我文件移除；其中自我文件移除僅排除 \texttt{candidate\_id=query\_id} 的同一篇判決書。第二類是 cut-off 搜尋：系統不一定固定保留前 5 名，而是於 validation 上以 grid search 比較三種策略。 \texttt{fixed\_topk} 固定保留前 $k$ 名；\texttt{largest\_gap\_cutoff} 在前 $N$ 名中尋找最大分數落差作為截斷點；\texttt{ratio\_cutoff} 則保留分數高於第一名一定比例的候選。各策略的最佳參數與結果見第~7 節。
